\subsection{Stochastic Dynamical Systems} 
\vspace{-2mm}
Consider a set of random vectors $\Statee_0, \ldots, \Statee_\horizon \in \states$ indexed at times $0, \ldots, \horizon$ and with state space $\states\subseteq\mathbb{R}^n$. A realization of this stochastic process is a sequence of values $\statee_1, \ldots, \statee_\horizon$, denoted as system trajectory  $\trajreal_{\statee_0}$. The joint distribution over $\Statee_1, \ldots, \Statee_\horizon$ is the trajectory distribution $\dist$, while the marginal distribution of $\Statee_0$ is known as the initial state distribution $\distinit$. It is assumed that $\distinit$ has support over a compact set of initial states $\init$, implying $\Pr[\statee_0 \notin \init] = 0$.

\mypara{Training and Deployment Environments} In the training environment, we pre-record or simulate datasets to conduct reachability analysis. Conversely, the deployment environment refers to the real world where we apply our reachable sets. There is typically a difference between the distribution of trajectories in the training and deployment environments. We refer to this difference as distribution shift. In this paper, we assume that for a predefined distribution on initial states $\statee_0 \sim \mathcal{W}$, the real-world trajectories $\trajreal_{\statee_0}$ are sampled from $\trajreal_{\statee_0} \sim \dist$, whereas the simulated trajectories $\trajsim_{\statee_0}$ are sampled from $\trajsim_{\statee_0} \sim \distzero$. 

\vspace{-2mm}
\subsection{Surrogate Model: Reachability \& Error Analysis}
 \vspace{-2mm}
 A surrogate model $\overallf: \init \times \Theta \to \states^\horizon$, with trainable parameters $\theta \in \Theta$, can be trained by sampling $\horizon$-step trajectories $\trajsim_{\statee_0} \sim \distzero$ from the simulator to predict the trajectory $\trajsim_{\statee_{0}} \in \states^\horizon$ given its initial state $\statee_0 \in \init$.  We call this dataset $\traindataset$, and we denote the predicted trajectory by,
\begin{equation}\label{eq:surrogate}
\bar{\traj}_{\statee_{0}} = \overallf(\statee_{0}\ ; \theta),\ \text{where},\ \  \overallf(\statee_{0}\ ; \theta) = 
\left[\mathsf{F}^1(\statee_0), \ldots, \mathsf{F}^n(\statee_0),  \ldots, 
\mathsf{F}^{(\horizon-1)n+1}(\statee_0), \ldots, \mathsf{F}^{n\horizon}(\statee_0)\right]^\top 
\end{equation}
\navid{where}, $\mathsf{F}^{(\timeid-1)n+\ell}(\statee_0)$ is the $\ell^{th}$ state component at the $\timeid^{th}$ time-step in the trajectory\footnote{Here, the dimension and time steps are stacked into a single vector.}. Let $e_\ell\in\reals^n$ denote the $\ell$-th basis vector of $\reals^n$. For a trajectory $\statee_1, \ldots \statee_\horizon$, and $\statee_0 \sim \mathcal{W}$, for $j =(\timeid-1)n+\ell$, we define the prediction errors as,
\begin{equation}
\label{eq:compwise_residual}
R^j =
e_{\ell}^\top \statee_\timeid - \mathsf{F}^j(\statee_0), \quad \timeid \in[\horizon], \ell\in[n].
\end{equation}


In this paper, we \navidd{introduce the residual $\rho:\mathbb{R}^{n\horizon} \to \mathbb{R}_{\ge 0}$ as a function} of the prediction errors $R^j$, where $j \in [n\horizon]$. In this section, we present a previously proposed example of such a \navid{function and later introduce an adjustment to reduce the conservatism in probabilistic reachability}.
\begin{definition}[Simulation \& Real Residual Distribution]
If the residual is generated by $\trajsim_{\statee_0}\sim \distzero$, we denote the residual distribution as $\ressim \sim \distzeroR$ where $\distzeroR$ is the simulation residual distribution. Conversely, if the residual is generated by $\trajreal_{\statee_0}\sim \dist$, we denote it as $\resreal \sim \distR$ where $\distR$ is the real residual distribution.
\end{definition}
We also utilize total variation, \cite{takezawa2005introduction} as a metric to quantify their distribution shift, $\tau\geq 0$. In other word, $\tau = \tv(\distzeroR, \distR)$, where $\tv$ refers to the total variation.

\mypara{Surrogate Flowpipe and Star-Set}  The surrogate flowpipe $\bar{X}\subset \reals^{n\horizon}$ is defined as a superset of the image of $\overallf(\init\ ;\theta)$. Formally, for all  $\statee_0\in \init$, we need that $\overallf(\statee_0\ ;\theta) \in \bar{X}$. Due to the recent achievements in verifying neural networks with ReLU activations, we limit ourselves to the choice of $\relu$ neural networks as our surrogate models, and we rely on the NNV toolbox from \cite{tran2020nnv} to compute the surrogate flowpipe.  Although other activation functions can be used, we anticipate more conservative results if non-ReLU activation functions are utilized. \navidg{We choose to use NNV in our analysis here, because it can yield accurate reachability results in settings where neural networks with $\relu$ activation function are used. } The approach in \cite{tran2020nnv} employs star-sets (an extension of zonotopes) to represent the reachable set and utilizes two main methods: (1) the exact-star method, which performs precise but slow computations, and (2) the approx-star method, which is faster but it is more conservative.
\begin{definition}[Star set \cite{bak2017simulation}]\label{def:star} A \textit{star set} $Y\subset \mathbb{R}^d$ is a tuple $\langle c, V, P \rangle$ where $c \in \mathbb{R}^d$ is the center, $V = \{v_1, v_2, \ldots, v_m \}$ is a set of $m$ vectors in $\mathbb{R}^d$ called \textit{basis vectors}, and $P : \mathbb{R}^m \rightarrow \{\top, \bot\}$ is a predicate. The basis vectors are arranged to form the star's $d \times m$ basis matrix. \navid{Given variables $\mu_\ell \in \mathbb{R}, \ell = 1,\ldots,m$}, the set of states represented by the star is given as:
\begin{equation}
 Y  = \left\{ y \mid y = c + \sum_{\ell=1}^{m} (\mu_\ell v_\ell) \text{ s.t. } P(\mu_1, \ldots, \mu_m) = \top \right\}.     
\end{equation}
\end{definition}

% \mypara{Quantification of the Distribution Shift} We use Total Variation (TV), \cite{takezawa2005introduction} as a metric to estimate the distribution shift. In this context, we do not have direct access to the distributions $\dist$ and $\distzero$. Instead, we are limited to a set of pre-recorded data points. Furthermore, computing or even estimating the joint distribution is impractical as it does not scale with the dimensionality of the multivariate random vectors like trajectories. However, since the residual is a measurable scalar function of the trajectories, the data processing inequality \cite{beaudry2011intuitive} ensures that $TV(\distzeroR, \distR) < TV(\dist, \distzero)$. Thus, by considering the residuals instead of the trajectories, we can efficiently compute for a lower bound on $TV(\dist, \distzero)$. 

\vspace{-2mm}
\subsection{Conformal Inference \& Probabilistic Reachability }
\vspace{-2mm}
\navid{A key step toward probabilistic reachability is to provide a provable $\delta$-quantile for the residual}. Let $\ressim_1 < \ressim_2 < \ldots < \ressim_L$ represent $L$ different i.i.d. residuals sampled from $\distzeroR$, and sorted \navid{in ascending order}. Given a confidence probability,  $\delta \in (0,1) $, a provable tight upper bound for the $\delta$-quantile of the residuals $\resreal \sim \distR$ is computable from samples $\ressim_i\sim \distzeroR , i\in[L]$, using robust conformal inference proposed in \cite{cauchois2020robust} that is an extension of CI, proposed in \cite{vovk2012conditional}.

\navid{The theory of conformal inference  states that for a new sample \(\ressim \sim \distzeroR\), the rank \(\ell := \lceil (L+1)\delta\rceil \le L\) satisfies $\Pr[ \ressim < \ressim_\ell] \geq \delta.$ This implies that \(\ressim_\ell\) serves as a provable upper bound for the \(\delta\)-quantile of \(\distzeroR\). However, this result does not extend to another residual \(\resreal \sim \distR\), which is drawn from a different distribution. To address this distribution shift, the theory of robust conformal inference introduces an adjustment to conformal inference. It establishes that for any random variable \(\resreal \sim \distR\) satisfying \(\tv(\distR, \distzeroR) \leq \tau\) with a threshold \(\tau > 0\), we have $\Pr[\resreal < \ressim_{\ell^*}] > \delta$, where  
\begin{equation}\label{eq:robustrank}
\ell^*:=\lceil (L+1)(1+1/L)(\delta+\tau)\rceil,\ \ \ell^* \leq L.
\end{equation}  
Thus, \(\ressim_{\ell^*}\) serves as an upper bound for the \(\delta\)-quantile of \(\distR\).}

% The theory of conformal inference \cite{vovk2012conditional} \navid{shows that,} for a new sample $\ressim \sim \distzeroR$, considering the rank $\ell:=\lceil (L+1)\delta\rceil\le L$, we have,
% $\Pr[ \ressim < \ressim_\ell] \geq \delta$. This implies $\ressim_\ell$ serves as provable upper-bound for the $\delta$-quantile of $\distzeroR$. However, this result does not hold for another residual, $\resreal \sim \distR$ as it is sampled from a different distribution. In this case, \navid{the theory of robust conformal inference \cite{cauchois2020robust} proposes an adjustment on conformal inference to consider distribution shift. Robust conformal inference shows that, given a threshold $\tau>0$, for any i.i.d random variable $\resreal \sim \distR$ where, $\tv(\distR, \distzeroR) \leq \tau$, we have, $\Pr[\resreal\! < \!\ressim_{\ell^*}]\! >\! \delta$, and,
% \begin{equation}\label{eq:robustrank}
% \ell^*=\lceil (L+1)(1+1/L)(\delta+\tau)\rceil.
% \end{equation}
% Thus, $\ressim_{\ell^*}$ is an upper-bound for the $\delta$-quantile of $\distR$.}

\mypara{Inflating Hypercube} In reachability analysis, the main purpose for the definition of the residual is to achieve a \navid{bounding region} that will cover the random sequence of prediction errors, $\PE = [R^1, R^2, \ldots, R^{n\horizon}]$ with a  confidence $\delta\in (0,1)$. In this case, as suggested by \cite{cleaveland2023conformal}, the $\max()$ operator over the absolute value of all errors is a suitable choice. The authors in \cite{hashemi2024statistical}, for some positive constants $\alpha_j, j\in[n\horizon]$ (details on the choice of $\alpha_j$ can  be found in \cite{hashemi2024statistical}), define   the residual
\begin{equation}\label{eq:Rmax}
\rho:= R = \max( \alpha_1 |R^1|, \alpha_2 |R^2|, \ldots, \alpha_{n\horizon}|R^{n\horizon}|),
\end{equation}
and \navidd{show} that such a bounding region is achievable by computing an upper bound for the $\delta$-quantile of residual, $R$. In other words, assuming $R^*$ as the mentioned upper bound, we have,
% \begin{equation*}%\label{eq:conclusion}
\begin{equation}\label{eq:Pstar}
\Pr[R<R^*] \geq \delta \!\! \navid{\iff} \!\!\Pr[P^*(R^1,\ldots,R^{n\horizon})\! = \!\top] \geq \delta,\ \   P^*(R^1,\ldots,R^{n\horizon})\!\! = \!\!\! \bigwedge_{j=1}^{n\horizon}\! (|R^j| \! < \! \frac{R^*}{\alpha_j})
\end{equation}
where $R^*$ is efficiently obtainable via robust conformal inference. \navidd{As defined in \eqref{eq:Pstar}, the predicate $P^*$ implies} that for every component $R^j, j=1,\ldots,n\horizon$ of the vector $\PE$ we have $-R^*/\alpha_j\leq R^j \leq R^*/\alpha_j$. This describes a hypercube which can be formulated as the following star set:
\begin{equation}
\delta X = \langle\  0_{n\horizon \times 1},\  I_{n\horizon},\  P^*(R^1,\ldots,R^{n\horizon}) \ \rangle \subset \mathbb{R}^{n\horizon}.
\end{equation}
Since $\Pr[P^* = \top] \geq \delta$, this star set serves as such a bounding region for $\PE$. We refer to this bounding region as inflating hypercube.


\mypara{$\delta$-Confident Flowpipe \& Probabilistic Reachability} For a given confidence probability $\delta\in (0,\ 1)$, and $\statee_0 \sim \mathcal{W}$, we say that $X \subseteq \reals^{n\horizon}$ is a $\delta$-confident flowpipe if for any random trajectory $\trajreal_{\statee_0} \sim \dist$, we have $\Pr[\trajreal_{\statee_0} \in X ] \geq \delta$. \navidd{In this paper, our ultimate goal is to propose a $\delta$-confident flowpipe. To compute such a flowpipe, the authors in \cite{hashemi2024statistical}} suggest simulating a set of trajectories, $\trajsim_{\statee_0} \sim \distzero, \statee_0 \sim \mathcal{W}$ and training a $\relu$ NN surrogate model $\overallf(\statee_0 ; \theta)$ on this dataset. This model will be utilized to compute for its surrogate reachset, $\bar{X}\subset \mathbb{R}^{n\horizon}$. They also sample a new set of trajectories $\trajsim_{\statee_0} \sim \distzero, \statee_0 \sim \mathcal{W}$ for error analysis on $\overallf(\statee_0 ; \theta)$ through robust conformal inference to compute for another hypercube $\delta X \subset \mathbb{R}^{n\horizon}$, known as the inflating hypercube, that covers the prediction errors $R^j , j\in[n\horizon]$ for trajectories $\trajreal_{\statee_0} \sim \dist$, with a provable probabilistic guarantee, and finally they propose the following lemma to compute for the $\delta$-confident flowpipe on $\trajreal_{\statee_0} \sim \dist, \statee_0 \sim \mathcal{W}$. See \cite{hashemi2024statistical} for the proof.
\begin{lemma}\label{lem:inclusion_conformal}
Let $\bar{X}$ be a surrogate flowpipe of the surrogate model $\overallf$ for the set of initial conditions $\init$. Let $\PEreal:=\left[ \errreal{1} , \errreal{2}, \ldots , \errreal{n\horizon} \right] $ be the sequence of prediction errors for $\trajreal_{\statee_0} \sim \dist$, where $\statee_0 \sim \mathcal{W}$, and  let $\delta X$ be the inflating hypercube for $\PEreal$ such that $\Pr[\PEreal \in \delta X] > \delta$. Then the inflated reachset $X = \bar{X} \oplus \delta X $ is a $\delta$-confident flowpipe for $\trajreal_{\statee_0} \sim \dist$ where $\statee_0 \sim \mathcal{W}$.
\end{lemma}

\vspace{-2mm}
\subsection{Problem Definition}
\vspace{-2mm}
We are interested in computing  a $\delta$-confident flowpipe $X$ from a set of trajectories $\trajsim_{\statee_0}$ collected from $ \distzero$ so that $X$ is also valid  for all trajectories $\trajreal_{\statee_0} \sim \dist$ when the total variation between $\distR$ and $\distzeroR$ is less than $\tau>0$. 
% However, this problem has been addressed before by authors in \cite{hashemi2024statistical}, and \textbf{our focus} here is to improve the scalability and the accuracy of this technique.
\navidd{While we are motivated by the results in \cite{hashemi2024statistical} which propose a  solution to the stated problem, we note that their solution lacks scalability and accuracy that results in sometimes large levels of conservatism, i.e., the set $X$ is unnecessarily large.}

\navid{The primary sources of conservatism and inaccuracy in the methodology described in \cite{hashemi2024statistical} stem from the training process for the surrogate model $\overallf(\statee_0\ ; \theta)$ and the method used to compute the inflating hypercube $\delta X$. In the following sections, we address both issues and propose solutions to improve the accuracy and scalability of this approach for the reachability analysis.}


% \mypara{$\delta$-Confident Flowpipe} For a given confidence probability $\delta\in (0,\ 1)$, a distribution $\ressim \sim \distzeroR$, $\statee_0 \sim \mathcal{W}$, and a radius $\tau\geq 0$, we say that $X \subseteq \reals^{n(\horizon+1)}$ is a $\delta$-confident flowpipe if for any random trajectory $\trajreal_{\statee_0} \sim \dist$, with $\tv(\distR, \distzeroR) \leq \tau$, we have $\Pr[\trajreal_{\statee_0} \in X ] \geq \delta$. 